% Figure: Hall-effect geometry. A current-carrying bar in a transverse magnetic
% field develops a transverse Hall voltage; the sign of that voltage gives the
% carrier type and its magnitude gives the carrier density. Plain TikZ.
\begin{figure}[htbp]
\centering
\begin{tikzpicture}[font=\scriptsize, >=Stealth]
% the conducting bar
\filldraw[engblue!12, draw=engblue, thick] (0,0) rectangle (6,2);
% current in along x
\draw[->, engred, very thick] (-1.4,1) -- (-0.1,1);
\node[engred, anchor=south] at (-0.8,1.05) {$I_x$};
\draw[->, engred, very thick] (6.1,1) -- (7.2,1);
\node[engred, anchor=south] at (6.7,1.05) {$I_x$};
% magnetic field out of plane (z)
\foreach \x in {1.2,2.4,3.6,4.8} {
  \foreach \y in {0.7,1.4} {
    \node[engblack] at (\x,\y) {$\odot$};
  }
}
\node[engblack, anchor=south] at (3,2.15) {$B_z$ out of page};
% Hall voltage probes across y
\draw[engorange, thick] (3,2) -- (3,2.6);
\node[engorange, circle, fill=engorange, inner sep=1.3pt] at (3,2.6) {};
\draw[engorange, thick] (3,0) -- (3,-0.6);
\node[engorange, circle, fill=engorange, inner sep=1.3pt] at (3,-0.6) {};
\node[engorange, anchor=west] at (3.15,2.6) {$+$};
\node[engorange, anchor=west] at (3.15,-0.6) {$-$};
\draw[<->, engorange] (4.7,-0.5) -- (4.7,2.5);
\node[engorange, anchor=west] at (4.75,1) {$V_H$};
% Lorentz deflection arrow inside the bar
\draw[->, enggray, thick] (1.5,1) -- (1.5,1.7);
\node[enggray, anchor=west, font=\tiny] at (1.55,1.4) {$\mathbf{F}=q\,\mathbf{v}\times\mathbf{B}$};
% width and thickness labels
\draw[<->, enggray] (-0.3,0) -- (-0.3,2);
\node[enggray, anchor=east] at (-0.35,1) {$w$};
\node[enggray, anchor=north, font=\tiny] at (3,-0.05) {thickness $t$ (into page)};
\end{tikzpicture}
\caption[Hall-effect measurement geometry]{A current $I_x$ flows along a bar of
width $w$ and thickness $t$ in a transverse magnetic field $B_z$. The magnetic
force deflects the moving carriers toward one face until the transverse electric
field it builds balances the deflection; the resulting transverse Hall voltage
is $V_H = I_x B_z/(n e t)$. The sign of $V_H$ identifies the carrier as electron
or hole, and its magnitude gives the carrier density $n$, the measurement worked
in the text. The same geometry, read in reverse with a known carrier density,
is a magnetic-field sensor.}
\label{fig:vol04ch12:hall-bar}
\end{figure}
